
\def\PUDcodigo{EME130}

\def\PUDcodacad{04506.56}

\def\PUDdisciplina{Projeto de Engenharia}

\def\PUDsemestre{Opt}

\def\PUDhoras{80}

%NOVOS ---------------------------------------------------
\def\PUDhorasTeoria{80}
\def\PUDhorasPratica{0}

\def\PUDinterECA{---}

\def\PUDatuaisECA{---}

\def\PUDrecursos{Projetor multimídia; Quadro branco e pincel.}

%----------------------------------------------------------

\def\PUDcreditos{4}

\def\PUDprerequisito{---}

\def\PUDpreacad{---}

\def\PUDobjetivos{Identificar as necessidades a serem abordadas no projeto;  Conhecer as fases de elaboração de um projeto;  Conhecer as ferramentas de elaboração de projetos;  Conhecer os modelos de planejamento de produtos industriais.}

\def\PUDementa{Introdução à Metodologia de Projeto em Engenharia;  Inovação Tecnológica em Produtos Industriais e Bens Materiais;  Desenvolvimento de Produtos;  Morfologia do Processo de Projeto.}

\def\PUDconteudo{ & Introdução à metodologia de projeto em engenharia: Processo de projeto.  Informações no projeto.  Viabilidade de produtos.  Tipos de produtos.  Requisitos de projeto.  Criatividade.  Análise do valor.  Projeto preliminar e Projeto detalhado.  Apresentação e Competição dos protótipos. 
\\ 
 & Inovação tecnológica em produtos industriais e bens materiais: Análises diacrônica e sincrônica dos modelos de planejamento de produto industrial.  Modelo de planejamento de produto industrial (PPI): projeto, produção e promoção.  Mercado, produção, desenho e sua integração.  Projetação no planejamento de produtos industriais. 
\\ 
 & Desenvolvimento de produtos: Identificação de problemas projetuais, técnicas analíticas projetuais, técnicas de geração e avaliação de alternativas, etapas do desenho do projeto, comunicação e especificações para a produção, realização de modelos (maquetes, mocapes, protótipos). 
\\ 
 & Morfologia do processo de projeto: Análise de informações e demanda.  Tipo de produtos e requisitos de projeto.  Síntese de soluções alternativas.  Função síntese.  Valoração e análise de valores.  Aspectos econômicos.  Projeto preliminar.  Seleção da solução.  Formulação dos modelos.  Materiais e processos de fabricação.  Projeto detalhado e revisão.  Atividades de Laboratório. }

\def\PUDmetodo{Aulas expositórias que podem ser teóricas e/ou práticas, onde as práticas no laboratório serão marcadas no decorrer da disciplina.}

\def\PUDavalia{A avaliação será feita com aplicação de provas teóricas e/ou práticas, além da possibilidade de inclusão de trabalhos e seminários no decorrer da disciplina.}

\def\PUDbibbasica{ &  \href{www.biblioteca.ifce.edu.br/index.asp?codigo_sophia=62218}{Fonseca}, José Wladimir Freitas da.  Elaboração e Análise de Projetos – A Viabilidade Econômico-Financeira.  São Paulo - SP.  Editora ATLAS.  2012.  209p.  ISBN: 9788522467518. 
\\ 
 &  \href{www.biblioteca.ifce.edu.br/index.asp?codigo_sophia=40465}{Pahl}, G.;  Beitz, W.;  Feldhusen, J.;  Grote, K.  Projeto na engenharia: fundamentos do desenvolvimento eficaz de produtos, métodos e aplicações.  São Paulo - SP.  Editora EDGARD BLUCHER.  2013.  412p.  ISBN: 9788521203636. 
\\
&  \href{www.biblioteca.ifce.edu.br/index.asp?codigo_sophia=62311}{Creswell}, John W.  Projeto de Pesquisa – Método Qualitativo, Quantitativo e Misto.  3º ed.  São Paulo - SP.  Editora Artmed.  2016.  296p.  ISBN: 9788536323008. }
%\\ 
 %&  POLAK, P.;  Projetos em Engenharia.  1ª edição.  São Paulo.  Editora HEMUS.  2005.  247p.  ISBN: 9788528905373.

\def\PUDbibcomplementar{ &  \href{www.biblioteca.ifce.edu.br/index.asp?codigo_sophia=61981}{Almeida}, Mário de Souza.  Elaboração de Projeto, TCC, Dissertação e Tese – Uma Abordagem Simples, Prática e Objetiva.  2º ed.  São Paulo - SP.  Editora Atlas.  2014.  82p.  ISBN: 9788522491155. 
\\ 
 &  \href{www.biblioteca.ifce.edu.br/index.asp?codigo_sophia=65145}{Madureira}, Omar Moore de.  Metodologia do projeto: planejamento, execução e gerenciamento.  2º ed.  São Paulo - SP.  Editora Blucher.  2015. 361p ISBN 9788521209133. 
\\
 &  \href{www.biblioteca.ifce.edu.br/index.asp?codigo_sophia=63128}{Juvinall}, Robert C.  Fundamentos do projeto de componentes de máquinas.  5º ed.  Rio de Janeiro - RJ.  LTC.  2016.  562.  ISBN: 9788521630098.
\\
 &  \href{www.biblioteca.ifce.edu.br/index.asp?codigo_sophia=58215}{Silva}, Ailton Santos.  Desenho técnico.  E-book.  Pearson.  136p.  ISBN: 9788543010977.
\\
 &  \href{www.biblioteca.ifce.edu.br/index.asp?codigo_sophia=58590}{Mott}, Robert L.  Elementos de máquina em projetos mecânicos. E-book. 5º ed.  Pearson.  924p.  ISBN: 9788543005904. }

\def\PUDautor{Venceslau Xavier de Lima Filho}

\def\PUDdata{2013-06-19}

\def\PUDrevisao{1}

\def\PUDrevisor{Samuel Vieira Dias}

\def\PUDdatarevisao{2017-04-30}

\PUDcorpo
